\begin{mistake}{2026-05-28}{03}

\begin{problemcontent}
设 \(f(u,v)\) 有连续偏导数，且 \(f(0,0)=0\)，判断下列等式中成立的是：
\[
\begin{array}{ll}
\text{(A)} & f(x,y)=x\displaystyle\int_0^1 \frac{\partial f(tx,y)}{\partial u}\,dt,\\[8pt]
\text{(B)} & f(x,y)=y\displaystyle\int_0^1 \frac{\partial f(x,ty)}{\partial v}\,dt,\\[8pt]
\text{(C)} & f(x,y)=x\displaystyle\int_0^1 \frac{\partial f(tx,ty)}{\partial u}\,dt+y\displaystyle\int_0^1 \frac{\partial f(tx,ty)}{\partial v}\,dt,\\[8pt]
\text{(D)} & f(x,y)=x\displaystyle\int_0^1 \frac{\partial f(tx,y)}{\partial u}\,dt+y\displaystyle\int_0^1 \frac{\partial f(x,ty)}{\partial v}\,dt.
\end{array}
\]
\end{problemcontent}

\begin{solutioncontent}
取从 \((0,0)\) 到 \((x,y)\) 的直线路径
\[
(u,v)=(tx,ty),\qquad 0\leq t\leq 1.
\]
令
\[
g(t)=f(tx,ty),
\]
则
\[
g(0)=f(0,0),\qquad g(1)=f(x,y).
\]
由二元复合函数链式法则，
\[
\begin{aligned}
g'(t)
&=\frac{\partial f(tx,ty)}{\partial u}\cdot \frac{d(tx)}{dt}
+\frac{\partial f(tx,ty)}{\partial v}\cdot \frac{d(ty)}{dt}\\
&=x\frac{\partial f(tx,ty)}{\partial u}
+y\frac{\partial f(tx,ty)}{\partial v}.
\end{aligned}
\]
由微积分基本公式，
\[
\begin{aligned}
f(x,y)-f(0,0)
&=g(1)-g(0)\\
&=\int_0^1 g'(t)\,dt\\
&=x\int_0^1 \frac{\partial f(tx,ty)}{\partial u}\,dt
+y\int_0^1 \frac{\partial f(tx,ty)}{\partial v}\,dt.
\end{aligned}
\]
又 \(f(0,0)=0\)，故
\[
f(x,y)=x\int_0^1 \frac{\partial f(tx,ty)}{\partial u}\,dt
+y\int_0^1 \frac{\partial f(tx,ty)}{\partial v}\,dt.
\]
因此正确选项为 \(\mathrm{C}\)。

检验其他选项：
\[
x\int_0^1 \frac{\partial f(tx,y)}{\partial u}\,dt=f(x,y)-f(0,y),
\]
故 A 需额外满足 \(f(0,y)=0\)；
\[
y\int_0^1 \frac{\partial f(x,ty)}{\partial v}\,dt=f(x,y)-f(x,0),
\]
故 B 需额外满足 \(f(x,0)=0\)。D 的右边一般为
\[
2f(x,y)-f(0,y)-f(x,0),
\]
不一定等于 \(f(x,y)\)。
\end{solutioncontent}

\mistakeknowledge{
\begin{itemize}
  \item \textbf{核心公式/定理}：若 \(g(t)=f(u(t),v(t))\)，则 \(g'(t)=f_u(u(t),v(t))u'(t)+f_v(u(t),v(t))v'(t)\)；并用 \(g(1)-g(0)=\int_0^1 g'(t)\,dt\)。
  \item \textbf{易错点}：A、B 对应的路径起点分别是 \((0,y)\)、\((x,0)\)，不是 \((0,0)\)，不能直接利用 \(f(0,0)=0\)。
  \item \textbf{关联知识}：判断此类积分恒等式时，要先看路径 \((u(t),v(t))\) 的起点和终点，再用链式法则还原函数增量。
\end{itemize}}

\end{mistake}
